\documentclass[10pt,a4paper]{amsart}
\usepackage[margin=19mm]{geometry}
\usepackage{amsmath,amssymb,mathtools}
\usepackage[scr=rsfs,cal=euler]{mathalfa}
\usepackage{xfrac}
\usepackage[unicode,hidelinks,bookmarks=false]{hyperref}
\newcommand{\ie}{i.e.\ }
\newcommand{\cf}{cf.\ }
\newcommand{\resp}{respectively\ }
\newcommand{\Subsection}{\subsection}
\providecommand{\nobreakdash}{\nobreak\hbox{-}}
\setlength{\emergencystretch}{2em}
\multlinegap0pt
\usepackage{amssymb}
\usepackage{mathtools}
\usepackage[shortlabels]{enumitem}
\usepackage{enumerate}
\usepackage{tikz}
\usepackage{xcolor}
\newtheorem{lemma}{Lemma}
\title{Eigenvalue growth of the discrete Hodge Laplacian across dimensions}
\author{Philipp Bartmann, Matthias Keller}
\date{Source: arXiv:2608.11170v2 (CC BY 4.0)}
\begin{document}
\maketitle

\noindent\emph{Source and licence.} Eigenvalue growth of the discrete Hodge Laplacian across dimensions. Version arXiv:2608.11170v2 (CC BY 4.0), distributed by arXiv under the Creative Commons Attribution 4.0 International licence, http://creativecommons.org/licenses/by/4.0/, as recorded in the arXiv licence metadata for this exact version and shown on the abstract page https://arxiv.org/abs/2608.11170v2. Full licence notice: \texttt{licenses/arxiv-2608.11170-CC-BY-4.0.txt}.\par\smallskip

\noindent\textit{Editorial scope.} Counted unit: the article's lemma claims (source0.tex lines 152--156). The article's complete original proof of it is delivered with it (source0.tex lines 157--175, one \texttt{proof} environment of 19 lines). No mathematics is written, repaired or re-derived: the statement and the proof are the article's own lines, in the article's own notation, and no retained line is substituted or re-flowed.  No further source-local context is needed: every cross-reference of every retained line resolves inside the two delivered spans.  Nothing was omitted from any retained span, so the package carries no omission record.  No retained line contains a citation, so the wrapper carries no bibliography and no bibliography entry.  No retained line contains a figure environment, an image-inclusion command or drawing code, so no image is delivered and the package declares no asset dependency.  Editorial formatting only, in addition: an AMS \texttt{amsart} preamble that restates the article's own declarations and macros used by the retained lines, namely the environments \texttt{lemma} on separate counters, together with the standard packages those lines need.  The retained environments print in this document as Lemma 1 in order of appearance, whereas the article prints the same items with the numbering of its own full document.  The complete native source of the article as distributed in the arXiv e-print for this exact version is bundled unchanged as \texttt{sources/207396/source0.tex}.
\begin{lemma}\label{claims}
    Let $\phi\in\ell^2(\Sigma_k),~k\geq 0.$ Then,
    \[(k+1)\Vert\phi\Vert^2 = \sum_{v\in V}\Vert\phi_v\Vert^2\quad\] 
    and for $k\geq 1,$ \[\quad k\Vert\partial_{k-1}\phi\Vert^2 = \sum_{v\in V}\Vert\partial_{k-2}\phi_v\Vert^2.\]
\end{lemma}

\begin{proof}
    The first identity follows as 
    \begin{align*}
        \sum_{v\in V}\Vert\phi_v\Vert^2 = \sum_{v\in V}\sum_{\substack{\rho\in\Sigma_{k-1}\\v\rho\in\Sigma_k}}\vert\phi(v\rho)\vert^2 = (k+1)\sum_{\tau\in\Sigma_k}\vert\phi(\tau)\vert^2 = (k+1)\Vert\phi\Vert^2.
    \end{align*}
    The factor $(k+1)$ arises, because for every $\tau\in\Sigma_k$ there are exactly $k+1$ tuples $(v,\rho)\in V\times\Sigma_{k-1}$ such that $\tau = v\rho.$
    The second equality follows in a similar way: For any $\tau' = v\rho$ with $v\in V$, $\rho\in\Sigma_{k-2}$ and any $\tau\succ\rho$ such that $v\tau\in\Sigma_k$, one finds that 
    \begin{equation}\label{thetas}
        0 = \delta_{k-1}\delta_{k-2} 1_\rho(v\tau) = \theta(\rho,\tau')\theta(\tau',v\tau)+\theta(\rho,\tau)\theta(\tau,v\tau).
    \end{equation}
    Thus,
    \begin{align*}
        \sum_{v\in V}\Vert\partial_{k-2}\phi_v\Vert^2 &=\sum_{v\in V} \sum_{\rho\in\Sigma_{k-2}}\Big\vert\sum_{\substack{\tau\succ\rho\\v\tau\in\Sigma_k}}\theta(\rho,\tau)\theta(\tau,v\tau)\phi(v\tau)\Big\vert^2\\
        & = \sum_{v\in V}\sum_{\rho\in\Sigma_{k-2}}\Big\vert\sum_{\substack{\tau\succ\rho\\v\tau\in\Sigma_k}}\theta(\rho,v\rho)\theta(v\rho,v\tau)\phi(v\tau)\Big\vert^2\\
        & = k\sum_{\tau'\in\Sigma_{k-1}}\Big\vert\sum_{\sigma\succ\tau'}\theta(\tau',\sigma)\phi(\sigma)\Big\vert^2\\
        &=k\Vert\partial_{k-1}\phi\Vert^2.
    \end{align*}
In the second equality, we used (\ref{thetas}) and in the third equality, we identified $\tau' = v\rho$ and $\sigma = v\tau.$
\end{proof}

\end{document}
